\SourceLeaf{98}%
imagin\"arer Gr\"ossen sich durch eine unendliche Ebene sinnlich machen, worin jeder Punkt, durch Abscisse $=a$, Ordinate $=b$ bestimmt, die Gr\"osse $a+bi$ gleichsam repr\"asentirt. Der stetige \"Ubergang von einem Werthe von $x$ zu einem andern $a+bi$ geschieht demnach durch eine Linie und ist mithin auf unendlich viele Arten m\"oglich. Ich behaupte nun, dass das Integral $\int \varphi x\mathord{.} dx$ nach zweien verschiednen \"Uberg\"angen immer einerlei Werth erhalte, wenn innerhalb des zwischen beiden die \"Uberg\"ange repr\"asentirenden Linien eingeschlossenen Fl\"achenraumes nirgends $\varphi x = \infty$ wird. Dies ist ein sehr sch\"oner Lehrsatz\HistoricalNote{\textsuperscript{*})}{*)}{Eigentlich ist hiebei noch angenommen, dass $\varphi x$ selbst eine einf\"ormige Function von $x$ ist, oder wenigstens f\"ur deren Werthe innerhalb jenes ganzen Fl\"achenraumes nur Ein System von Werten ohne Unterbrechung der Stetigkeit angenommen wird.}, dessen eben nicht schweren Beweis ich bei einer schicklichen Gelegenheit geben werde. Er h\"angt mit sch\"onen andern Wahrheiten, die Entwicklungen in Reihen betreffend, zusammen. Der \"Ubergang nach jedem Punkte l\"asst sich immer ausf\"uhren, ohne jemals eine solche Stelle wo $\varphi x=\infty$ wird zu ber\"uhren. Ich verlange daher, dass man solchen Punkten ausweichen soll, wo offenbar der urspr\"ungliche Grundbegriff von $\int \varphi x\mathord{.} dx$ seine Klarheit verliert und leicht auf Widerspr\"uche f\"uhrt. \"Ubrigens ist zugleich hieraus klar, wie eine durch $\int \varphi x\mathord{.} dx$ erzeugte Function f\"ur einerlei Werthe von $x$ mehrere Werthe haben kann, indem man nemlich beim \"Ubergange dahin um einen solchen Punkt wo $\varphi x=\infty$ entweder gar nicht, oder einmal, oder mehreremale herumgehen kann. Definirt man z.\,B. $\log x$ durch $\int \frac{1}{x}\,dx$, von $x=1$ anzufangen, so kommt man zu $\log x$ entweder ohne den Punkt $x=0$ einzuschliessen oder durch ein- oder mehrmaliges Umgehen desselben; jedesmal kommt dann die Constante $+2\pi i$ oder $-2\pi i$ hinzu: so sind die vielfachen Logarithmen von jeder Zahl ganz klar. Kann $\varphi x$ nie f\"ur einen endlichen Werth von $x$ unendlich werden, so ist das Integral immer nur eine einf\"ormige Function. Diess ist z.\,B. der Fall f\"ur $\varphi x = \dfrac{e^x-1}{x}$, so dass $\int \dfrac{e^x-1}{x}\,dx$ gewiss eine einf\"ormige Function von $x$ ist, deren Werth durch die \textls{immer} convergirende, also immer einen und nur \textls{Einen} Sinn habende Reihe dargestellt wird
\[
x+\tfrac{1}{4}xx+\tfrac{1}{18}x^3+\tfrac{1}{96}x^4+\text{etc.}
\]
\noindent Ich wollte, Herr \textsc{Soldner} h\"atte, da er doch einmal eine neue Function ein%
